%%%%%%%%%%  scan idx 288  |  folio 275  %%%%%%%%%%
\[
 \coprod_{i\geq0} L[2ni]\simeq
 L\oplus\left(\coprod_{i\geq0}L[2ni]\right)[2n],
\]
et d'appliquer a).

\textbf{Lemme 1.3.} Soient \(C\) une catégorie triangulée, \(A\) une sous-catégorie épaisse \(([\mathrm V]\ \mathrm I\ 2.1)\). Les foncteurs d'inclusion et de passage au quotient définissent une suite exacte
\[
 k(A)\longrightarrow k(C)\longrightarrow k(C/A)\longrightarrow 0.
\]

\emph{Preuve.} Voir (SGA 5 VIII 3.1).

\textbf{Lemme 1.4.} Soient \(A\) une catégorie additive (resp. abélienne), et \(f\) une fonction additive (1.1) sur \(\Ob K^b(A)\) (resp. \(\Ob D^b(A)\)). Pour \(E\in\Ob K^b(A)\) (resp. \(D^b(A)\)) on a
\[
 f(E)=\sum(-1)^i f(E^i)
\]
(resp.
\[
 f(E)=\sum(-1)^i f(H^i(E))
\]).

\emph{Preuve.} Laissée en exercice au lecteur.

\textbf{1.5.} Soient \(C\) une catégorie \emph{abélienne} et \(A\) une sous-catégorie pleine de \(C\) stable par sommes finies. Une application \(f\) de \(\Ob A\) dans un groupe abélien \(G\) est dite \emph{additive} si l'on a \(\sum(-1)^i f(E^i)=0\) pour tout complexe acyclique \(E\) de \(C\), à degrés bornés et dont les composantes appartiennent à \(A\). Lorsque \(A\) est stable par noyau d'épimorphisme (I 1.2), il revient au même de dire qu'on a \(f(L)=f(L')+f(L'')\) pour toute suite exacte \(0\longrightarrow L'\longrightarrow L\longrightarrow L''\longrightarrow0\) de \(C\) dont les termes appartiennent à \(A\). Les applications additives de \(\Ob A\) dans \(G\) forment un groupe abélien dépendant fonctoriellement de \(G\). Notons \(K^{b,\emptyset}(A)\) la sous-catégorie épaisse de \(K^b(A)\) formée des complexes acycliques. On déduit aisément de (1.3) et (1.4) que l'application
\begin{equation}\tag{1.5.1}
 \cl_A:\Ob A\longrightarrow k\bigl(K^b(A)/K^{b,\emptyset}(A)\bigr)
\end{equation}
